% Fragment for inclusion in the main report.
% Requires amsmath, amssymb, amsthm and theorem/corollary environments.
\section{The exact information lost by fixed congruence weights}
\label{sec:periodic-obstruction}

The earlier reports establish a mean-square obstruction for the standard
character of an $S_3$ cubic field. The same obstruction holds for every
nonlinear irreducible character of every finite Galois group. More generally,
one can compute the optimal periodic approximation to any Frobenius class
function exactly. This is an application of character theory and the classical
Chebotarev theorem; it is a diagnostic for proposed extensions, not a claim
to solve another open conjecture.

Let $M/\mathbb Q$ be a finite Galois extension with group $G$. Fix an integer
$q\ge1$ and put
\[
 A_q=(\mathbb Z/q\mathbb Z)^\times,\qquad
 E_q=M\cap\mathbb Q(\zeta_q),\qquad
 H_q=\operatorname{Gal}(M/E_q).
\]
For $q=1$, interpret $A_q$ as the trivial group. Since $E_q/\mathbb Q$ is
abelian, $H_q$ is normal and contains the commutator subgroup $[G,G]$.
Let $S$ be any fixed finite set containing the primes ramifying in $M$ and
the primes dividing $q$. At every $p\notin S$, choose a representative
$g_p\in G$ of the arithmetic Frobenius conjugacy class. If $f:G\to\mathbb C$ is a class
function, its value $f(g_p)$ is independent of this choice.

The restriction of $b\in A_q=\operatorname{Gal}(\mathbb Q(\zeta_q)/\mathbb Q)$
to $E_q$ determines one coset of $H_q$ in $G$. Define
\begin{equation}\label{eq:periodic-best}
 m_q(b)=\frac1{|H_q|}
   \sum_{\substack{g\in G\\g|_{E_q}=b|_{E_q}}}f(g).
\end{equation}
Thus $m_q$ is a finite, explicitly specified function on the reduced residue
classes. Write
\[
 \|f\|_G^2=\frac1{|G|}\sum_{g\in G}|f(g)|^2,
 \qquad
 \|a\|_{A_q}^2=\frac1{\varphi(q)}\sum_{b\in A_q}|a(b)|^2.
\]

\begin{theorem}[Optimal approximation by a fixed modulus]
\label{thm:periodic-projection}
For every class function $f:G\to\mathbb C$ and every function
$a:A_q\to\mathbb C$, both of the following limits exist and have the same
value:
\begin{align}
 \lim_{X\to\infty}\frac1{\pi(X)}
   \sum_{\substack{p\le X\\p\notin S}}
       |f(g_p)-a(p\bmod q)|^2
 &\;=\;\mathcal E_q(f,a),\label{eq:unweighted-mse}\\
 \lim_{X\to\infty}\frac1X
   \sum_{\substack{p\le X\\p\notin S}}
       (\log p)|f(g_p)-a(p\bmod q)|^2
 &\;=\;\mathcal E_q(f,a),\label{eq:weighted-mse}
 \intertext{where}
 \mathcal E_q(f,a)
 &=\|f\|_G^2-\|m_q\|_{A_q}^2
                  +\|a-m_q\|_{A_q}^2.\label{eq:mse-decomposition}
\end{align}
Consequently $m_q$ is the unique optimal approximation on $A_q$, and the
minimum error is $\|f\|_G^2-\|m_q\|_{A_q}^2$.

Let $f=\sum_{\chi\in\operatorname{Irr}(G)}c_\chi\chi$ be its expansion in
irreducible characters, with the usual normalized character inner product.
Then the minimum has the equivalent expression
\begin{equation}\label{eq:mse-character-energy}
 \min_a\mathcal E_q(f,a)
 =\sum_{\substack{\chi\in\operatorname{Irr}(G)\\
                    H_q\not\subseteq\ker\chi}}|c_\chi|^2.
\end{equation}
Here $H_q\subseteq\ker\chi$ means that the representation affording
$\chi$ is trivial on $H_q$.
\end{theorem}

\begin{proof}
Put $N=M\mathbb Q(\zeta_q)$. Restriction identifies its Galois group with
the fiber product
\[
 \Gamma=\{(g,b)\in G\times A_q:
                  g|_{E_q}=b|_{E_q}\}.
\]
Its projection onto $A_q$ is surjective and has fibers of size $|H_q|$,
so $|\Gamma|=\varphi(q)|H_q|$. The function
$F(g,b)=|f(g)-a(b)|^2$ is constant on conjugacy classes of $\Gamma$:
$f$ is a class function and $A_q$ is abelian. Chebotarev's density theorem,
summed over the finitely many conjugacy classes, gives
\begin{equation}\label{eq:fiber-product-average}
 \lim_{X\to\infty}\frac1{\pi(X)}
    \sum_{\substack{p\le X\\p\notin S}}F(g_p,p\bmod q)
 =\frac1{\varphi(q)|H_q|}
    \sum_{b\in A_q}
    \sum_{g|_{E_q}=b|_{E_q}}|f(g)-a(b)|^2.
\end{equation}
This is the unconditional fixed-field Chebotarev theorem. No GRH, effective
error estimate, or uniformity in $M$ or $q$ is being used.

For a fixed $b$, the mean of $f$ over the indicated coset is $m_q(b)$.
Expanding the square and using the vanishing of the mean of $f-m_q(b)$
on that coset gives the orthogonal decomposition
\[
 \frac1{|H_q|}\sum_{g|_{E_q}=b|_{E_q}}|f(g)-a(b)|^2
 =\frac1{|H_q|}\sum_{g|_{E_q}=b|_{E_q}}|f(g)-m_q(b)|^2
      +|m_q(b)-a(b)|^2.
\]
Average over $b$. Every $g\in G$ is counted the same number of times,
because the restriction $A_q\to\operatorname{Gal}(E_q/\mathbb Q)$ is
surjective. This proves \eqref{eq:mse-decomposition} and uniqueness of the
minimum, since every residue class has the positive weight
$1/\varphi(q)$.

For completeness the logarithmic prime weight does not change the limit.
Write $U(X)$ for the unweighted numerator in
\eqref{eq:unweighted-mse}. Chebotarev and the ordinary prime number theorem
give
$U(X)=\mathcal E_q(f,a)\operatorname{Li}(X)+o(X/\log X)$.
Partial summation yields
\[
 \sum_{\substack{p\le X\\p\notin S}}(\log p)F(g_p,p\bmod q)
 =U(X)\log X-\int_2^X\frac{U(t)}t\,dt
 =\mathcal E_q(f,a)X+o(X).
\]
Changing the fixed lower endpoint alters only a constant. This proves
\eqref{eq:weighted-mse}. Equivalently one may normalize the two sums by
the number of retained primes and by their sum of logarithms, respectively:
those denominators are asymptotic to $\pi(X)$ and $X$. Deleting or adding
finitely many unramified primes does not affect either limit.

It remains to identify the character components retained by the coset
average. Let $\rho$ be an irreducible representation of $G$ with character
$\chi$, and set
\[
 P_H=\frac1{|H_q|}\sum_{h\in H_q}\rho(h).
\]
This is projection onto the subspace fixed by $H_q$. Because $H_q$ is
normal, that subspace is $G$-invariant; by irreducibility it is either
zero or the whole representation. Therefore
\[
 \frac1{|H_q|}\sum_{h\in H_q}\chi(gh)
 =\operatorname{tr}(\rho(g)P_H)
 =\begin{cases}
   \chi(g),&H_q\subseteq\ker\chi,\\
   0,&H_q\not\subseteq\ker\chi.
  \end{cases}
\]
The characters in the first case factor through the abelian quotient
$G/H_q$ and hence have degree one. Irreducible-character orthogonality
now makes the squared norm lost by this projection exactly the sum in
\eqref{eq:mse-character-energy}.
\end{proof}

\begin{corollary}[Every nonlinear irreducible trace is invisible to fixed congruences]
\label{cor:nonlinear-periodic}
If $\chi$ is an irreducible character of $G$ of degree greater than one,
then for every fixed $q$ and every $a:A_q\to\mathbb C$,
\begin{equation}\label{eq:nonlinear-mse}
 \mathcal E_q(\chi,a)
 =1+\frac1{\varphi(q)}\sum_{b\in A_q}|a(b)|^2.
\end{equation}
In particular the unique optimal fixed-periodic approximation is zero,
and the minimum mean-square error is exactly one in both
\eqref{eq:unweighted-mse} and \eqref{eq:weighted-mse}.
\end{corollary}

\begin{proof}
Since $G/H_q$ is abelian, the nonlinear irreducible representation cannot
be trivial on $H_q$. The character average in the proof of
Theorem~\ref{thm:periodic-projection} is consequently zero on every coset,
so $m_q=0$. Character orthogonality gives $\|\chi\|_G^2=1$.
\end{proof}

The normalization here is the ordinary trace character. If the trace is
divided by the representation dimension $d$, its minimum mean-square
error is instead $1/d^2$.

\paragraph{What the best fixed modulus can retain.}
Let $M^{\mathrm{ab}}=M^{[G,G]}$ be the maximal abelian subextension.
By Kronecker--Weber, it is contained in $\mathbb Q(\zeta_Q)$ for some
integer $Q$. For every multiple $q$ of $Q$ we have
$E_q=M^{\mathrm{ab}}$ and $H_q=[G,G]$. Indeed every abelian subextension
of $M$ lies in $M^{\mathrm{ab}}$, while the assumed divisibility gives
the reverse containment for $E_q$. Thus for such $q$,
\begin{equation}\label{eq:all-fixed-periodic}
 \min_a\mathcal E_q(f,a)
   =\sum_{\substack{\chi\in\operatorname{Irr}(G)\\\chi(1)>1}}
       |c_\chi|^2.
\end{equation}
The right side is also the infimum over all fixed moduli and all periodic
approximations. An arbitrary class function is exactly recoverable from
some fixed residue modulus in prime mean square if and only if it factors
through the abelianization of $G$. For smaller moduli one can additionally
lose one-dimensional character components whose abelian subextensions
are not contained in the cyclotomic field chosen.

\paragraph{The cubic example.}
For an irreducible cubic polynomial with Galois group $S_3$, the trace of
the standard two-dimensional representation is the number of roots
modulo an unramified prime minus one. Its values and frequencies are
\[
\begin{array}{c|ccc}
\text{Frobenius type}&\text{identity}&\text{transposition}&\text{three-cycle}\\
\hline
\chi(g)&2&0&-1\\
\text{prime density}&1/6&1/2&1/3
\end{array}
\]
and its mean square is $4/6+1/3=1$. Even after the sign of Frobenius is
known, the conditional mean trace is zero: on $A_3$ it is
$(2-1-1)/3=0$, and on the odd coset every trace is zero. The general
theorem proves that no finer fixed congruence system improves this
approximation. It applies both to $X^3-X-1$ and to $X^3-2$, whose
particular instances were treated in the earlier reports.

\paragraph{Implication for a proposed transfer.}
Every function of a fixed residue class is a finite linear combination of
Dirichlet characters, by Fourier expansion on $A_q$. Consequently a
construction whose new prime weight is obtained solely by a fixed finite
combination of Dirichlet-character weights cannot reproduce a nonlinear
irreducible Frobenius trace, even in the mean-square sense above. A
successful extension must use some additional arithmetic information or
a different limiting mechanism.

The quantifier \emph{fixed} is essential. The theorem first fixes $q$ and
$a$ and then lets $X$ tend to infinity. It makes no uniform assertion
when $q=q(X)$ grows or when the residue weight is adaptively chosen from
the primes up to $X$. In fact, if $q>X$ is prime, each prime $p\le X$
occupies its own residue class, so an artificially chosen function can
interpolate all of those finitely many traces exactly. Such interpolation
supplies neither useful analytic control nor an arithmetic construction,
but it shows why one must not advertise \eqref{eq:nonlinear-mse} as an
obstruction to every possible use of growing moduli. Nor does the result
prove logical independence of Artin or Hecke GRH from Dirichlet GRH.
